\section{Literature Survey}
\label{sec:litreview}

We surveyed recent work on hands-on security education, adaptive learning, machine
learning for education, and gamification. We asked three questions of every source:
\emph{why} it matters, the \emph{tools} or platform it used, and its
\emph{drawbacks} or the gap it leaves open. Table~\ref{tab:litsurvey} summarises the
survey along those three questions, and the section closes by stating how our work
improves on this body of work.

\begin{small}
\renewcommand{\arraystretch}{1.25}
\begin{longtable}{@{}p{1.7cm} p{4.5cm} p{2.6cm} p{4.6cm}@{}}
\caption{Literature survey: why, tools, and drawbacks.}
\label{tab:litsurvey}\\
\toprule
\textbf{Study} & \textbf{Why (focus \& finding)} & \textbf{Tools / platform} & \textbf{Drawbacks / gap left open} \\
\midrule
\endfirsthead
\multicolumn{4}{@{}l}{\small\itshape Table~\ref{tab:litsurvey} continued}\\
\toprule
\textbf{Study} & \textbf{Why (focus \& finding)} & \textbf{Tools / platform} & \textbf{Drawbacks / gap left open} \\
\midrule
\endhead
\midrule
\multicolumn{4}{r@{}}{\small\itshape continued on next page}\\
\endfoot
\bottomrule
\endlastfoot

\cite{Svabensky2020what} & Systematic review of a decade of security-education papers. Most measure perceptions, not behaviour. & Literature review. & Effectiveness rarely measured by skill or behaviour. \\

\cite{Chapman2014picoctf} & Game-based CTF competition motivates novices and scales widely. & PicoCTF. & Fixed difficulty, outcome-only, no personalisation. \\

\cite{Electronics2022teaching} & Closest design point: gamifies OWASP Top~10 web material with micro-learning and CTF. & Custom CTF + LMS. & No adaptive difficulty, no ML, drawbacks of gamification not addressed. \\

\cite{Streicher2016adaptive} & Framework for personalised and adaptive serious games (learner model + adaptation). & Adaptive game engines. & General framework, not built for a web-security CTF class. \\

\cite{Vykopal2022smart} & Adaptive training matches task difficulty to the learner (114 students). & KYPO / adaptive engine. & Infrastructure-heavy, not web-app / browser focused. \\

\cite{Svabensky2024detecting} & Predicts at-risk students from logged actions using ML (313 students). & ML on action logs. & Detects \emph{that} a student struggles, not a tailored response, and not web-specific. \\

\cite{Toda2018dark} & Names the negative effects of gamification. Leaderboards most linked to harm. & Review. & Warns of harm, offers few concrete design remedies. \\

\cite{Almeida2023negative} & Systematic mapping of 87 papers: points, badges, competition can demotivate. & Systematic mapping. & Confirms the problem, solutions left to designers. \\

\cite{Karagiannis2020analysis} & Evaluates open-source CTF platforms. CTFd is the most teaching-oriented. & CTFd and others. & Limited scaffolding, no adaptive difficulty, weak assessment. \\

\cite{Owasp_juiceshop,Ctfd_docs} & Mature, reusable targets and scoring: a vulnerable app and a CTF framework. & Juice Shop, CTFd. & Standalone, not integrated with adaptivity or teacher analytics. \\

\end{longtable}
\end{small}

\subsection{Hands-on and gamified security education works}

Practical work teaches security well, and CTF is the main vehicle for it. A
systematic review of a decade of security-education papers frames the field and
notes that most studies measure what learners \emph{think} rather than what they
\emph{do}~\cite{Svabensky2020what}. CTF as a teaching tool traces back to game-based
competitions such as PicoCTF~\cite{Chapman2014picoctf}, and the published design
closest to ours gamifies OWASP Top~10 web material through micro-learning and CTF
mechanics~\cite{Electronics2022teaching}. The lesson is that hands-on, gamified
teaching works, so the question is not \emph{whether} to build it but \emph{how} to
fix what existing designs leave out.

\subsection{Every learner is different: adaptive learning and personalisation}

A single fixed sequence of challenges suits nobody exactly. A framework for
personalised and adaptive serious games pairs a learner model with rules that change
the game to suit that learner~\cite{Streicher2016adaptive}, and an adaptive
cybersecurity-training study with 114 students showed that matching task difficulty
to the learner produces a better fit than a fixed path~\cite{Vykopal2022smart}. The
gap is that these ideas are rarely built into a web-security CTF platform that
can be run for a class.

\subsection{Predicting who will struggle with machine learning}

Personalisation needs a signal: some estimate of how each learner is doing. The most
relevant work predicts at-risk students from their logged actions across two
learning environments and 313 students, and argues that because instructor time is
limited, automated detection of struggling students is
needed~\cite{Svabensky2024detecting}. A \emph{simple, interpretable} model over
features the platform already records (attempts, accuracy, time on task, hint use)
can flag who needs help and feed an adaptive recommender.

\subsection{Gamification and its documented drawbacks}

Two reviews document the harm directly. The element most often linked to harm is the
leaderboard, and the most-cited harm is loss of motivation and performance for
weaker learners~\cite{Toda2018dark}. A systematic mapping of 87 papers confirms that
points, badges, competition, and leaderboards are the mechanics most often blamed
for negative effects~\cite{Almeida2023negative}. This is the direct reason our
gamification module uses a peer-banded leaderboard (ranks near the viewer rather
than one absolute ranking) and keeps competition optional, designing gamification
\emph{against} these documented failure modes rather than ignoring them.

\subsection{Platforms, vulnerable applications, and tools}

The tools already exist and should be reused. An evaluation of open-source CTF
platforms finds CTFd the most teaching-oriented and a sound reference for scoring,
hints, and challenge dependencies~\cite{Karagiannis2020analysis,Ctfd_docs}. For the
vulnerable application, OWASP Juice Shop is a mature target covering the OWASP
Top~10~\cite{Owasp_juiceshop}, whose challenge families are fixed by the OWASP
Top~10 itself~\cite{Owasp2021top10}. What none of these tools provides is a single
classroom platform combining a vulnerable app, automated scoring, drawback-aware
gamification, ML-driven personalisation, and an instructor dashboard.

\subsection{How our work improves on existing systems (gaps)}
\label{sec:gaps}

The survey distils into six gaps, each addressed by reusing mature tools and
keeping any machine learning simple. Each gap states how our work improves on prior systems, and
each is answered by exactly one objective (Section~\ref{sec:objectives}).

\begin{table}[htbp]
\centering
\caption{Gaps (how our work improves on existing systems), traced to evidence and objective.}
\label{tab:gaps}
\renewcommand{\arraystretch}{1.3}
\begin{tabularx}{\linewidth}{@{}c X p{2.7cm} c@{}}
\toprule
\textbf{ID} & \textbf{Gap / how our work improves} & \textbf{Evidence} & \textbf{Obj.} \\
\midrule
G1 & Vulnerable apps are standalone targets. We integrate one into a single platform
     that also scores, gamifies, and analyses.
   & \cite{Owasp_juiceshop,Karagiannis2020analysis} & O1 \\
G2 & Game elements are bolted on without regard to harm. We design them against the
     documented negative effects.
   & \cite{Toda2018dark,Almeida2023negative} & O2 \\
G3 & Assessment is flag-only or manual. We score on completion, accuracy, and time,
     producing a signal a model can learn from.
   & \cite{Chapman2014picoctf,Karagiannis2020analysis,Svabensky2024detecting} & O3 \\
G4 & The path is one-size-fits-all. We add adaptive hints and a personalised next
     challenge, because every learner is different.
   & \cite{Streicher2016adaptive,Vykopal2022smart} & O4 \\
G5 & Instructor visibility stops at solve counts. We predict who is at risk and show
     it on a dashboard so support arrives early.
   & \cite{Svabensky2024detecting} & O5 \\
G6 & Effectiveness is under-evaluated. We measure engagement and practical skill
     gain in a pilot.
   & \cite{Svabensky2020what} & O6 \\
\bottomrule
\end{tabularx}
\end{table}
